Výroková logika pracuje s tzv. \emph{výroky} -- intuitivně vzato větami, o nichž
můžeme prohlásit, zda jsou pravdivé, či lživé. Výroky pojíme dohromady tzv.
\emph{logickými spojkami}, kterých je obecně mnoho, ale my se soustředíme na ty
v matematice snad nejběžněji užívané.

\begin{definition}[Výrok]
  \label{def:vyrok}
  Výrokem nazveme jakoukoli větu, o níž můžeme prohlásit, zda platí, či nikoliv.
\end{definition}
\lessboxspace
\begin{warning}
  Na skutečnost, zda je nějaká věta výrokem, nebo ne, nemá žádný vliv to, jestli
  jsme schopni \emph{rozhodnout} o její platnosti. I věty, o nichž nevíme, zda
  jsou pravdivé, či ne, ale víme, že buď jsou, nebo nejsou, jsou výroky.
\end{warning}
\begin{example}
  \label{exam:vyrok}
  Věty
  \begin{itemize}
    \item Zítra bude pršet.
    \item Mám rád kočky.
    \item Čenda je druhý příchod Ježíše.
  \end{itemize}
  \textbf{jsou} výroky, zatímco věty
  \begin{itemize}
    \item Přineš mi kafe!
    \item Co bych si tak mohl zahrát.
    \item Pochopím tohle někdy?
  \end{itemize}
  výroky \textbf{nejsou}.
\end{example}

Pravda a lež, jsouce základními logickými koncepty, mají své vlastní nezávislé
jedinečné symboly, třeba $T$ (pro \uv{true}) a $F$ (pro \uv{false}). Poslední
sadou symbolů výrokové logiky jsou symbol negace ($\neg$) a logické spojky ($
\wedge $, $ \vee $, $ \Rightarrow $ a $ \Leftrightarrow $). Povíme si, co
znamenají.

\begin{definition}[Negace]
  \label{def:negace}
  Ať $p$ je nějaký výrok. Potom $\neg p$ (čteno \uv{ne $p$} nebo \uv{negace
  $p$}) je výrok s opačnou pravdivostní hodnotou od $p$, čili, jest-li $p$
  pravda, pak $\neg p$ je lež a naopak.
\end{definition}

\begin{example}
  \label{exam:negace}
  Pokud za $p$ označíme výrok: \uv{Chutná mi pivo,} pak $\neg p$ je výrok:
  \uv{Nechutná mi pivo.}
\end{example}

Logické spojky $ \wedge $, $ \vee $, $ \Rightarrow $ a $ \Leftrightarrow $
slučují dva výroky do jednoho, jehož pravdivostní hodnota (tj. fakt, zda je
pravdivý, či lživý) závisí na hodnotě slučovaných výroků. Představíme si
\emph{konjunkci} a \emph{disjunkci}, jež odpovídají českým spojkám \emph{a} a
\emph{nebo}.

\begin{definition}[Konjunkce]
  \label{def:konjunkce}
  Ať $p$ a $q$ jsou výroky. \emph{Konjunkce} (či lidštěji logické \emph{a}) je
  logická spojka se symbolem $ \wedge $ s vlastností, že $p \wedge q$ (čten
  \uv{$p$ a $q$}) je pravda \textbf{jedině tehdy}, když $p$ je pravda i $q$ je
  pravda.
\end{definition}

\begin{example}
  \label{exam:konjunkce}
  Je-li $p$ výrok: \uv{Venku sněží,} a $q$ je výrok: \uv{Je mi zima,} pak $p
  \clr{ \wedge } q$ je výrok: \uv{Venku sněží \textbf{\clr{a}} je mi zima.}
\end{example}

\begin{definition}[Disjunkce]
  \label{def:disjunkce}
  At́ $p$ a $q$ jsou výroky. \emph{Disjunkce} (či lidštěji logické \emph{nebo})
  je logická spojka se symbolem $ \vee $ s vlastností, že $p \vee q$ (čten
  \uv{$p$ nebo $q$}) je pravda, \textbf{jedině tehdy}, když aspoň jeden (ale lze
  i oba) z výroků $p$ a $q$ je pravda.
\end{definition}

\begin{example}
  \label{exam:disjunkce}
  Je-li $p$ výrok: \uv{Neumím počítat,} a $q$ je výrok: \uv{Jsem prostě mimo,}
  pak $p \clr{ \vee } q$ je výrok: \uv{Neumím počítat
  \textbf{\clr{nebo}} jsem prostě mimo}.
\end{example}

Poslední dva symboly -- $ \Rightarrow $ a $ \Leftrightarrow $ -- patří implikaci
a ekvivalenci, což jsou logické spojky užívané k vyjádření podmínky.

\begin{definition}[Implikace]
  \label{def:implikace}
  Ať $p$ a $q$ jsou výroky. \emph{Implikace} je logická spojka se symbolem $
  \Rightarrow $ s vlastností, že $p \Rightarrow q$ (čten třeba \uv{Když $p$, tak
  $q$,} nebo \uv{Z $p$ plyne $q$.}) je pravda \textbf{v těchto dvou případech}:
  \begin{itemize}
    \item $p$ je pravda a $q$ je pravda;
    \item $p$ je lež.
  \end{itemize}
  Výroku $p$ často říkáme \emph{předpoklad} a výroku $q$ \emph{závěr}.
\end{definition}
\begin{example}
  \label{exam:implikace}
  Jsou-li $p$ a $q$ výroky po řadě: \uv{Vstanu včas,} a \uv{Stihnu bus,} tak $p
  \clr{ \Rightarrow } q$ je výrok: \uv{\textbf{\clr{Když}} vstanu včas,
  \textbf{\clr{tak}} stihnu bus.}
\end{example}
\begin{warning}
  Logická definice implikace $ \Rightarrow $ je v lehké oposici jejímu
  přirozenému vnímání v češtině. Totiž, je-li předpoklad $p$ lež, pak celý výrok
  $p \Rightarrow q$ je pravda. Matematická logika se v tomto případě řídí
  pravidlem, že \textbf{ze lži může plynout jakýkoli závěr}.
\end{warning}
\lessboxspace
\begin{definition}[Ekvivalence]
  \label{def:ekvivalence}
  \emph{Ekvivalence} je logická spojka se symbolem $ \Leftrightarrow $
  s~vlastností, že $p \Leftrightarrow q$ (čten \uv{$p$ právě tehdy, když $q$})
  je pravda, když výroky $p$ a $q$ mají stejnou pravdivostní hodnotu, tj.
  \textbf{v těchto dvou případech}:
  \begin{itemize}
    \item $p$ je pravda a $q$ je
      pravda;
    \item $p$ je lež a $q$ je lež.
  \end{itemize} Je-li ekvivalence $p
  \Leftrightarrow q$ pravdivá, výroky $p$ a $q$ nazýváme \emph{ekvivalentní}.
\end{definition}
\begin{example}
  \label{exam:ekvivalence}
  Je-li $p$ výrok: \uv{Čekám deset minut na přestup,} a $q$ výrok: \uv{Spěchám
  do školy,}, pak $p \clr{ \Leftrightarrow } q$ je výrok: \uv{Čekám deset minut
  na přestup \textbf{\clr{právě tehdy, když}} spěchám do školy.}
\end{example}
